%% =========================================================================
\section{Exercise 1: Synchronous Machine Control and Initialization}
%% =========================================================================

%% ── SIMULINK MODEL DESCRIPTION ────────────────────────────────────────────
\subsection{Simulink Model}

\figref{fig:ex1_diagram} shows the complete Simulink model used for Exercise~1.
The model comprises a \MVA{555} round-rotor synchronous machine connected to a
resistive-inductive load via a three-phase bus. The AVR (AC1C exciter) regulates
the terminal voltage, and the governor--turbine subsystem controls the
mechanical power input. Two relay-controlled breakers implement the load step
events at $t = \SI{3}{s}$ (step off) and $t = \SI{9}{s}$ (step on).
Measurement blocks capture $V_t$, $\omega$, $P$, $Q$, and $I_f$ for
post-processing.

\begin{figure}[H]
  \centering
  \includegraphics[width=0.75\textwidth]{EX1_figure.png}
  \caption{Exercise 1 Simulink model: single synchronous machine with AVR,
    governor, and switched loads.}
  \label{fig:ex1_diagram}
\end{figure}

%% ── TASK 1 ────────────────────────────────────────────────────────────────
\subsection{Task 1 -- Properties and Functionality of Simulation Components}

\subsubsection{Synchronous Machine}

The synchronous machine is the main component in the model. It converts
mechanical shaft power into balanced three-phase electrical power. It is a
round-rotor type, which is standard for high-speed steam-turbine generators in
thermal power stations. The model uses the Park (d-q~axis) transformation to
project the three-phase quantities onto a rotating two-axis reference frame.
This gets rid of the time-varying inductances and makes the equations easier
to work with.

The machine parameters are listed in \tabref{tab:machine_params}
(values taken from~\cite{ieee1110}).

\begin{table}[H]
  \centering
  \caption{Synchronous Machine Parameters}
  \label{tab:machine_params}
  \begin{tabular}{>{\bfseries}l l l}
    \toprule
    \textbf{Parameter} & \textbf{Symbol} & \textbf{Value} \\
    \midrule
    Rated apparent power       & $S$          & \MVA{555}   \\
    Rated voltage (L-L RMS)    & $V_\text{rated}$ & \kV{24} \\
    Rated frequency            & $f_0$        & \Hz{60}     \\
    Pole pairs                 & $p$          & 1           \\
    Armature resistance        & $R_a$        & 0.003\,pu   \\
    d-axis synchronous reactance & $X_d$      & 1.81\,pu    \\
    q-axis synchronous reactance & $X_q$      & 1.76\,pu    \\
    d-axis transient reactance   & $X_d'$     & 0.30\,pu    \\
    q-axis transient reactance   & $X_q'$     & 0.65\,pu    \\
    d-axis subtransient reactance & $X_d''$   & 0.23\,pu    \\
    q-axis subtransient reactance & $X_q''$   & 0.25\,pu    \\
    Inertia constant           & $H$          & 3.525\,s    \\
    Damping coefficient        & $D$          & 0.01\,pu    \\
    \bottomrule
  \end{tabular}
\end{table}

The steady-state active and reactive power delivered to the load are
approximated by the power-angle equations, derived from the phasor diagram of
the synchronous machine:

\begin{equation}
  P \approx \frac{E \cdot V}{X} \sin\delta
  \label{eq:active_power}
\end{equation}

\begin{equation}
  Q \approx \frac{V}{X}\left(E\cos\delta - V\right)
  \label{eq:reactive_power}
\end{equation}

\noindent where $E$ is the internal EMF (pu), $V$ is the terminal voltage (pu),
$X$ is the synchronous reactance (pu), and $\delta$ is the power angle
(electrical degrees). From \eqnref{eq:active_power}, active power is controlled
mainly through the power angle $\delta$, which is set by the prime mover.
\eqnref{eq:reactive_power} shows that reactive power depends on the internal
voltage $E$, set by field excitation. Because of this, the governor (P-f) and
AVR (Q-V) control loops can operate independently.

\subsubsection{AVR and Exciter (AC1C Model)}

The Automatic Voltage Regulator (AVR) and exciter constitute a closed-loop
feedback control subsystem maintaining the generator terminal voltage at the
reference value ($V_\text{ref} = \pu{1.0}$). The exciter model is the
AC1C~type~\cite{ieee4215}, which represents an alternator-rectifier excitation
system with field-current-based feedback.

\begin{figure}[H]
  \centering
  \includegraphics[width=0.65\textwidth]{AVR_exiter.png}
  \caption{AVR and exciter (AC1C) block diagram showing the closed-loop
    voltage regulation subsystem.}
  \label{fig:avr_exciter}
\end{figure}

The AVR measures the terminal voltage $V_t$ and compares it to $V_\text{ref}$.
The difference produces an error signal that adjusts the field voltage
$E_\text{fd}$ through the exciter. The simplified transfer function is:

\begin{equation}
  E_\text{fd}(s) = \frac{K_A}{1 + s T_A}\left(V_\text{ref} - V_t(s)\right)
  \label{eq:avr}
\end{equation}

\noindent where $K_A$ is the AVR gain and $T_A$ is the AVR time constant.
If $V_t$ drops below $V_\text{ref}$ (for example after a large load is
connected), the AVR increases $E_\text{fd}$. This raises the internal
EMF~$E$. From \eqnref{eq:reactive_power}, a higher~$E$ at roughly
constant~$\delta$ means more reactive power~$Q$ is injected, which pushes the
terminal voltage back up toward the reference. This is how voltage regulation
works in power systems.

\subsubsection{Governor and Prime Mover}

The governor and prime mover subsystem regulates rotor speed and system
frequency by controlling the mechanical torque input from the turbine. The
governor implements a speed droop (regulation) characteristic:

\begin{equation}
  R = \frac{\Delta f / f_\text{rated}}{\Delta P / P_\text{rated}} = 0.05
  \quad (5\%\text{ droop})
  \label{eq:droop_def}
\end{equation}

Model parameters are: speed reference $\omega_\text{ref} = \pu{0.998}$,
turbine time constant $T_t = \SI{0.2}{s}$, and governor time constant
$T_g = \SI{0.3}{s}$. The steady-state power output as a function of
frequency is:

\begin{equation}
  P = \frac{f_{nl} - f_\text{sys}}{S_P}
  \label{eq:droop_power}
\end{equation}

\noindent where $f_{nl}$ is the no-load (setpoint) frequency (pu), $f_\text{sys}$
is the system frequency (pu), and $S_P$ is the droop slope. A 5\% droop means
a 5\% reduction in frequency commands a 100\% increase in rated power output.

\subsubsection{Machine Inertia Block}

The machine inertia block represents the kinetic energy stored in the rotating
mass of the generator rotor and turbine shaft. This is characterised by the
inertia constant $H$:

\begin{equation}
  H = \frac{\tfrac{1}{2}J\omega_s^2}{S_\text{rated}}
  \quad [\text{seconds}]
  \label{eq:inertia_constant}
\end{equation}

\noindent For this machine, $H = 3.525$\,s. This means the rotor stores enough
kinetic energy to supply rated power for 3.525 seconds at synchronous speed.
The rotor dynamics follow the swing equation:

\begin{equation}
  \frac{2H}{\omega_s}\frac{d\omega}{dt} = P_m - P_e - D(\omega - \omega_s)
  \label{eq:swing}
\end{equation}

\noindent where $\omega$ is the rotor angular velocity (rad/s), $\omega_s$ is
synchronous speed (rad/s), $P_m$ is the mechanical power input (pu), $P_e$ is
the electrical power output (pu), and $D = 0.01$\,pu is the damping
coefficient. The swing equation is the main tool used in transient stability
analysis. The Rate of Change of Frequency (ROCOF) right after a
$\SI{100}{MW}$ load change is:

\begin{equation}
  \left.\frac{df}{dt}\right|_{t=0}
  = -\frac{f_s}{2H}\cdot\frac{\Delta P}{S_\text{rated}}
  = -\frac{60}{7.05}\cdot\frac{100}{555}
  \approx -\SI{1.54}{Hz/s}
  \label{eq:rocof}
\end{equation}

%% ── TASK 2 ────────────────────────────────────────────────────────────────
\subsection{Task 2 -- Variation of Synchronous Machine Parameters Under Load}

The simulation incorporates two load-switching events: a load disconnection
at $t = \SI{3}{s}$ (`Step off') and a larger load connection at $t = \SI{9}{s}$
(`Step on'). The system response is examined through five key variables.

\subsubsection{Active Power}

\begin{figure}[H]
  \centering
  \includegraphics[width=0.65\textwidth]{active_power.png}
  \caption{Active mechanical power $P_m$ variation under load switching events.
    Initial steady state: \MW{250}. At $t=\SI{3}{s}$: step down to \MW{150}.
    At $t=\SI{9}{s}$: step up to \MW{350}.}
  \label{fig:ex1_pm}
\end{figure}

The machine begins at steady-state delivering \MW{250} to the initial load.
At $t = \SI{3}{s}$, the `Step off' breaker disconnects a load block, causing
the electrical power demand to fall sharply. The governor responds by reducing
the turbine valve opening and the active power settles near \MW{150} after a
brief transient. At $t = \SI{9}{s}$, the `Step on' breaker connects a large
load block; demand jumps to approximately \MW{350}, creating a sudden
mechanical-to-electrical power deficit which the governor corrects by opening
the turbine valves wider.

\subsubsection{Rotor Speed and System Frequency}

\begin{figure}[H]
  \centering
  \includegraphics[width=0.65\textwidth]{angular_velocity.png}
  \caption{Rotor angular velocity $\omega$ (rad/s) variation under load switching
    events. Since the machine has one pole pair ($p=1$), electrical frequency
    equals mechanical speed in per-unit: $f_\text{pu} = \omega_\text{pu}$.}
  \label{fig:ex1_speed}
\end{figure}

Because the machine has a single pole pair ($p=1$), electrical frequency equals
mechanical speed in per-unit: $f_\text{pu} = \omega_\text{pu}$.

\textbf{At $t=\SI{3}{s}$ (load removal):} Electrical power $P_e$ falls below
mechanical power $P_m$. The right-hand side of \eqnref{eq:swing} becomes
positive, causing rotor acceleration. Speed peaks at approximately \pu{1.015}
before the governor reduces turbine output, and the system settles near
\pu{1.010}. The droop equation predicts the steady-state deviation:

\begin{equation}
  \Delta f_\text{ss} = -R \cdot \frac{\Delta P}{S_\text{rated}}
  = -0.05 \times \frac{100 \times 10^6}{555 \times 10^6}
  \approx -0.009\,\text{pu}
  \label{eq:droop_calc_3s}
\end{equation}

\textbf{At $t=\SI{9}{s}$ (load addition):} A large power deficit ($P_e \gg P_m$)
decelerates the rotor sharply to approximately \pu{0.980}. The governor
increases turbine power and frequency recovers to approximately \pu{0.990}
($\approx \SI{59.4}{Hz}$), consistent with \eqnref{eq:droop_calc_3s} applied
to the \MW{200} step increase.

\subsubsection{Terminal Voltage}

\begin{figure}[H]
  \centering
  \includegraphics[width=0.65\textwidth]{terminal_voltage.png}
  \caption{Terminal voltage $V_t$ (pu) showing AVR response. Voltage nadir
    at $t=\SI{9}{s}$: approximately \pu{0.82}. AVR restores $V_t$ to
    \pu{1.0} within 3--4 seconds.}
  \label{fig:ex1_vt}
\end{figure}

The terminal voltage holds at \pu{1.0} in steady state. At $t=\SI{3}{s}$, the
load reduction causes a minor transient overshoot ($\approx\pu{1.02}$) as the
AVR momentarily over-excites before damping. At $t=\SI{9}{s}$, the large load
addition drives a severe voltage dip to approximately \pu{0.82}, caused by the
increased armature current producing large internal voltage drops across $X_d$.
The AVR responds immediately by commanding maximum field excitation, and $V_t$
is restored to \pu{1.0} within 3--4 seconds, demonstrating the critical role of
fast AVR action in maintaining power quality during severe disturbances.

\subsubsection{Reactive Power}

\begin{figure}[H]
  \centering
  \includegraphics[width=0.65\textwidth]{Reacrtive_power.png}
  \caption{Reactive power $Q_m$ (MVAr) variation. Initial steady state:
    \MVAr{15}. Large transient surge at $t=\SI{9}{s}$ corresponds to AVR
    driving increased field excitation to restore $V_t$.}
  \label{fig:ex1_qm}
\end{figure}

Reactive power sits at approximately \MVAr{15} at the initial steady state.
Sharp transient spikes occur at $t=\SI{3}{s}$ and $t=\SI{9}{s}$, corresponding
to sudden changes in network impedance and the subsequent excitation adjustments
by the AVR. At $t=\SI{9}{s}$, the reactive power surges significantly before
settling, as the AVR drives a large increase in $E_\text{fd}$ to compensate for
the voltage dip, injecting additional reactive power into the network per
\eqnref{eq:reactive_power}.

\subsubsection{Angular Torque}

\begin{figure}[H]
  \centering
  \includegraphics[width=0.65\textwidth]{angular_torque.png}
  \caption{Angular torque variation during load switching events.
    Torque transients at $t=\SI{3}{s}$ and $t=\SI{9}{s}$ reflect the
    mechanical--electrical power imbalance described by the swing equation.}
  \label{fig:ex1_torque}
\end{figure}

The angular torque exhibits sharp transients at each switching event. At
$t=\SI{3}{s}$, the sudden reduction in electrical load causes a momentary
surplus of mechanical torque, accelerating the rotor. At $t=\SI{9}{s}$, the
large load addition creates a torque deficit, decelerating the rotor. The
governor adjusts the prime mover torque to restore equilibrium in both cases.

\subsubsection{Power Dissipated}

\begin{figure}[H]
  \centering
  \includegraphics[width=0.65\textwidth]{power_dissipated.png}
  \caption{Power dissipated in the load during switching events. The step
    changes at $t=\SI{3}{s}$ and $t=\SI{9}{s}$ correspond to the load
    breaker operations.}
  \label{fig:ex1_pdiss}
\end{figure}

The power dissipated in the load tracks the breaker switching events directly.
At $t=\SI{3}{s}$, the load disconnection reduces the dissipated power, and at
$t=\SI{9}{s}$, the additional load connection increases it. The transient
oscillations in dissipated power reflect the coupled electromechanical dynamics
of the synchronous machine as it adjusts to the new operating point.

\subsubsection{Relationships Between Key Variables}

\textbf{Active Power $\leftrightarrow$ Frequency:}
The swing equation~\eqnref{eq:swing} and governor droop~\eqnref{eq:droop_power}
link these two variables. When mechanical power exceeds electrical power
($P_m > P_e$), the rotor accelerates and frequency goes up. The opposite
happens when there is a deficit. The simulation matches this --- frequency rises
after the $t=\SI{3}{s}$ load removal and drops after the $t=\SI{9}{s}$ load
addition, following the active power response with a delay set by the machine
inertia. This P--f relationship is the basis of primary frequency regulation.

\textbf{Reactive Power $\leftrightarrow$ Terminal Voltage:}
Equations~\eqnref{eq:reactive_power} and \eqnref{eq:avr} link these through
the AVR. When terminal voltage drops, the AVR increases excitation
$\rightarrow$ $E$ rises $\rightarrow$ $Q$ output increases
$\rightarrow$ voltage is restored. The simulation shows this clearly: the
$t=\SI{9}{s}$ voltage dip comes with a sharp reactive power surge as the AVR
responds. In steady state, the Q--V loop and the P--f loop are mostly
independent, which allows voltage and frequency to be controlled separately.
